%-------------------------------------------------------------------------------
\section{Compatible Discrete Operator (CDO) schemes}

\subsection{Introduction }

Compatible Discrete Operator schemes have been recently introduced in \CS to
handle distorted and/or polyhedral meshes in a robust way. CDO schemes belong to
the broad family of \emph{compatible} or \emph{mimetic} or
\emph{structure-preserving} schemes. For more details about compatible
discretizations and previous works on this topic, please refer
to~\cite[Chapter~2 - Compatible spatial discretization]{Bonel14}. Constitutive
relations or closure laws (\emph{metric relations}) are clearly separated from
relations defining differential operators or conservation laws
(\emph{topological relations}).

CDO schemes gather several space discretization schemes according to the
localization of the degrees of freedom (DoFs). Namely, DoFs are defined and
located following their physical meaning (potential, circulation, flux or
density). DoFs and discrete differential operators are linked together to
satisfy a discrete De Rham complex mimicking the continuous one.


Figure~\ref{fig:cdo_dofs} depicts a sketch of all DoF locations available in CDO
schemes. According to the space discretization scheme, some of them are
considered to build a given CDO scheme.
\begin{itemize}
\item \emph{vertex-based} schemes when DoFs are located at mesh vertices. DoF
  definition corresponds to a discrete field associated with a potential.
\item \emph{vertex+cell-based} schemes when DoFs are both located at mesh cells
  and mesh vertices. DoF definitions correspond to a discrete field associated
  with a potential. After building the cellwise system, DoF associated with the
  cell mesh can be algebraically removed using an operation called \emph{static
    condensation}. After the resolution (a system of size equal to the number of
  mesh vertices for a scalar-valued equation), cell DoFs can be recovered
  performing the reverse operation.
\item \emph{face-based} schemes when DoFs are located at mesh faces and mesh
  cells. DoF definitions correspond to a discrete field associated with a
  potential. After building the cellwise system, DoF associated with the cell
  mesh can be algebraically removed using an operation called \emph{static
    condensation}. After the resolution (a system of size equal to the number of
  mesh faces for a scalar-valued equation), cell DoFs can be recovered
  performing the reverse operation.
\item \emph{edge-based} schemes when DoFs are located at mesh edges. DoF
  definition corresponds to the circulation along each edge of a vector-valued
  unknown. This scheme is only available for vector-valued \emph{curl-curl}
  problems.
\item \emph{cell-based} schemes when DoFs are located at mesh cells (for DoFs
  associated with a potential unknown) and mesh faces (for DoFs associated with
  a flux unknown). This scheme is only available for scalar-valued diffusion
  equations.
\end{itemize}

\begin{figure}[htbp]
  \centering
  \includegraphics[width=0.7\textwidth]{cdo_dofs}
  \caption{Location of the degrees of freedom (DoFs) on a hexagonal prismatic
    cell according to the CDO space discretization schemes.}
  \label{fig:cdo_dofs}
\end{figure}

\begin{remark}
  For a scalar-valued equation associated with a pure diffusion problem,
  \emph{vertex-based} and \emph{vertex+cell-based} schemes are related to the
  \emph{primal} formulation of the problem, \emph{cell-based} schemes are
  related to the \emph{mixed} formulation of the problem and \emph{face-based}
  schemes to a \emph{hybrid} formulation of the problem.
\end{remark}


CDO schemes have been devised and mathematically analyzed during three PhD
thesis. Part of these works have been integrated inside \CS and are now
available in \CS when activating the \texttt{CDO} mode:
\begin{itemize}
\item Bonelle's Phd~\cite{Bonel14} for vertex-based, face-based and cell-based
  schemes in the case of steady anisotropic diffusion equation and an
  edge-based scheme for a steady \emph{curl-curl} equation;
\item Cantin's Phd~\cite{Cant16} for the scalar-valued unsteady
  convection/diffusion/reaction equation with vertex-based schemes and
  vertex+cell-based schemes;
\item Milani's PhD~\cite{Mila20} for the discretization of the steady/unsteady
  Stokes and Navier--Stokes equations using face-based schemes. Two
  velocity-pressure coupling algorithms are investigated and analyzed: (1) a
  fully coupled (or \textit{monolithic}) algorithm and (2) an
  \textit{artificial compressibility} algorithm.
\end{itemize}

\subsection{Diffusion operator}

The discretization of the diffusion is detailed in several reports and
articles. The EDF's report~\cite{Bonel12} (in French) is a good
introduction. The link with other space discretizations and the mathematical
analysis (stability, convergence proof and \emph{a priori} error estimates) is
available in~\cite{BonEr14}. The key operator is the discrete Hodge
operator. This operator can be built in several ways. An example of the
influence of the choice of the discrete Hodge operator on the quality of the
solution is available in~\cite{BDPE15}.

An detailed explanation on the way to enforce Dirichlet boundary conditions in
vertex-based scheme along with the description of vertex+cell-based schemes for
scalar-valued diffusion equation is available in~\cite{Cant16}.

\subsubsection{Discrete Hodge Operator}

The discrete Hodge operator is the key operator linking the primal and dual
entities. It acts as a mass-like matrix mapping discrete differential forms
(degrees of freedom such as circulations or fluxes) associated with either the
primal or dual mesh. This operator bridges degrees of freedom (DoFs) in duality.

Depending on the discretization scheme and the localization of the DoFs, several
types of discrete Hodge operators are defined in \CS:
\begin{itemize}
\item \textbf{Primal Vertices to Dual Cells (VpCd):} Maps primal vertex values
  to dual cell values. This operator acts as a mass matrix for vertex-based
  schemes.
\item \textbf{Primal Edges to Dual Faces (EpFd):} Connects primal edge
  circulations (e.g., circulation of a vector field) to dual face fluxes. This
  operator plays an important role in the build of stiffness matrices for
  vertex-based schemes.
\item \textbf{Primal Faces to Dual Edges (FpEd):} Connects primal face fluxes to
  dual edge circulations. This operator plays an important role in cell-based
  schemes.
\item \textbf{Dual Edges to Primal Faces (EdFp):} Maps dual edge circulations to
  primal face fluxes. This operator plays an important role in the build of
  stiffness matrices for face-based schemes.
\item \textbf{Primal Cells to Dual Vertices (CpVd):} Maps primal cell values to
  dual vertex values. This operator acts as a mass matrix for cell-based
  schemes.
\item \textbf{Dual Vertices to Primal Cells (VdCp):} Maps dual vertex values to
  primal cell values.
\item \textbf{Face-based (FC):} Play the role of a mass-like matrix for
  face-based schemes. In this case, one operates on a hybrid space of degrees of
  freedom located at both faces and cells.
\item \textbf{Vertex+Cell-based (VC):} Play the role of a mass-like matrix for
  vertex+cell-based schemes, operating on a hybrid space of degrees of freedom
  situated at both vertices and cells.
\end{itemize}

\paragraph{Algorithms for discrete Hodge operators.}

To compute the local contributions of the discrete Hodge operators, \CS offers
four main algorithms, offering different trade-offs between mesh alignment,
consistency, and computational cost:

\begin{itemize}
\item \textbf{Voronoi algorithm:} This algorithm assumes that a geometric
  orthogonality condition holds between the primal and dual entities (e.g., the
  line connecting two cell centers is orthogonal to their shared face). It leads
  to a diagonal discrete Hodge operator (conceptually related to a Two-Point
  Flux Approximation - TPFA). While computationally very efficient, using this
  algorithm on non-orthogonal or skewed meshes introduces a consistency error
  that does not vanish under mesh refinement.
\item \textbf{Whitney on the Barycentric Subdivision (WBS):} This algorithm
  relies on a subdivision of each polyhedral cell into a set of sub-tetrahedra
  via its barycentric subdivision (see~\cite[Bonel14]). Classical Finite Element
  polynomial shape functions (Whitney forms) are then defined on this
  subdivision. This provides a consistent and robust approximation on arbitrary
  polyhedral meshes. On tetrahedral meshes, operators defined using this
  algorithm are similar to classical Finite Element formulations. This algorithm
  does not required a geometric orthogonality condition for the mesh.
\item \textbf{Orthogonal Consistency/Stabilization (COST):} This algorithm
  decomposes the local Hodge operator $\mathsf{H}_c$ into two orthogonal contributions for each cell $c\in C$ ($C$ is the set of primal cells):
  \begin{equation}
    \mathsf{H}_c = \mathsf{H}_c^{CO} + \beta^2 \mathsf{H}_c^{ST}
  \end{equation}
  where $\mathsf{H}_c^{CO}$ is the consistency part (ensuring exact
  reconstruction of piecewise constant vector-valued fields or piecewise linear
  scalar-valued fields) and $\mathsf{H}_c^{ST}$ is the stabilization part
  (guaranteeing stability and positive-definiteness on nearly arbitrary
  meshes). The user-defined scaling parameter $\beta^2$ controls the weight of
  the stabilization. This formulation underpins several advanced polyhedral
  schemes of the literature such as the SUSHI scheme (\textup{S}cheme
  \textup{U}sing \textup{S}tabilization and \textup{H}ybrid \textup{I}nterfaces)
  ($\beta = \frac{1}{\sqrt{3}}$), the Discrete Geometric Approach (DGA) scheme
  ($\beta = \frac{1}{3}$), and the Generalized Crouzeix--Raviart (GCR) scheme
  ($\beta = 1$). This algorithm also naturally supports anisotropic and
  tensor-valued material properties.
\item \textbf{Bubble stabilization:} Similar to the \textbf{COST} approach, this
  algorithm utilizes an orthogonal decomposition between consistency and
  stabilization. However, the stabilization term $\mathsf{H}_c^{ST}$ is built
  using piecewise linear bubble functions, similar to those of the Finite
  Element literature. This leads to a stabilization scaling of:
  \begin{equation}
    C_{stab} = 0.1 \times \beta^2
  \end{equation}
  where $\beta$ is the user-defined stabilization coefficient and $0.1$ is the
  scaling arising from the integral of the squared bubble function over a
  tetrahedral subdivision. This is particularly useful for vertex-based schemes
  to avoid over-stabilization on stretched cells.
\end{itemize}



\subsection{Convection operator}

Several advection schemes are detailed in~\cite{Cant16} for scalar-valued
vertex-based schemes. The classical upwind and centered discretizations are
discussed along with a \emph{P\'eclet-robust} scheme which automatically
applies an upwind weighting according to the value of the local P\'eclet
number.

An upwind and centered schemes are detailed in~\cite{Mila20} in the case of
scalar- and vector-valued face-based schemes.

\subsection{Stokes and Navier--Stokes equations}

The discretization of the Stokes or Navier--Stokes equations in the case of a
laminar flow with a constant mass density and viscosity is detailed
in~\cite{Mila20}. With these assumptions, the viscous term is simply a
vector-valued Laplacian operator.

Two types of velocity/pressure coupling algorithms are available:
\begin{itemize}
\item A \emph{monolithic} algorithm: a fully coupled algorithm between the
  velocity components and the pressure DoFs. This yields to a linear system
  which is a saddle-point problem. A specific algorithm to solve steady-state
  problem is also available.
\item An \emph{artificial compressibility} algorithm~\cite{GuMi15} up to a
  second-order accuracy~\cite{MiBoE22}. Only unsteady problems are tackled with this
  algorithm. This is an experimental feature.
\end{itemize}

A short introduction to this discretization describing the main operators in
the steady case with the monolithic approach is available in~\cite{BoErM20}.


%%% Local Variables:
%%% mode: LaTeX
%%% TeX-master: "theory"
%%% End:
